\documentclass[11pt]{article}
\usepackage[margin=1in]{geometry}
\usepackage{amsmath}
\usepackage{booktabs}
\usepackage{array}
\usepackage{hyperref}
\hypersetup{colorlinks=true,urlcolor=blue}

\setlength{\parindent}{0pt}
\setlength{\parskip}{0.6em}

\begin{document}

\begin{flushright}
Prakash Suryawanshi\\
Qodeigence Infocircle Private Limited\\
Maharashtra, India\\
\texttt{prakashsurywanshi1@live.com}
\end{flushright}

\vspace{0.5em}
\today

\vspace{1em}
The Editor-in-Chief\\
\textit{IEEE Internet of Things Journal}\\
% (Alternative target: IEEE Internet of Things Magazine --- adjust addressee accordingly.)
IEEE

\vspace{1em}
\textbf{Subject: Submission of an original research manuscript --- ``A Sustainable Hybrid QKD--PQC Security Architecture for Quantum-Resilient Residential IoT''}

\vspace{0.5em}
Dear Editor,

On behalf of my co-authors and myself, I am pleased to submit the original research
manuscript entitled ``A Sustainable Hybrid QKD--PQC Security Architecture for
Quantum-Resilient Residential IoT'' for consideration in the
\emph{IEEE Internet of Things Journal}. The manuscript has not been published
previously and is not under consideration for publication elsewhere.

\section*{Scope and Significance}
Battery-powered residential IoT nodes must remain operational for years while facing
the emerging threat of quantum cryptanalysis and the ``harvest-now, decrypt-later''
attack model. Migrating conventional public-key cryptography to standardized
post-quantum cryptography (PQC) directly on constrained leaf nodes, however, imposes a
severe energy and radio-bandwidth penalty because large lattice-based keys are split
across many IEEE~802.15.4/6LoWPAN frames. This paper addresses that sustainability
tension with a gateway-centric hybrid architecture that keeps asymmetric post-quantum
operations off the leaf nodes while preserving long-term post-quantum resilience. The
work fits the Journal's scope at the intersection of IoT security, edge/energy
sustainability, and next-generation network protocols.

\section*{Key Points and Contributions}
\begin{enumerate}
  \item \textbf{Sustainable hybrid QKD--PQC architecture.} An edge-centric Quantum-Safe
  Home Security Gateway (QHSG) structurally isolates security-intensive key-management
  operations from high-volume application data streams; last-mile Quantum Key
  Distribution (QKD) material is ingested where optical infrastructure exists, while
  standardized ML-KEM-768 (FIPS~203) provides a deployable path over conventional
  networks. Bulk data is protected by authenticated symmetric encryption.

  \item \textbf{Grounded security--sustainability model.} Endpoint compute energy, radio
  transmission overhead, and 6LoWPAN fragmentation penalties are formalized as explicit
  design parameters. Constants are anchored to the IEEE~802.15.4 physical layer and the
  TI~CC2652R datasheet ($P_{\mathrm{TX}}{=}22$~mW, $P_{\mathrm{RX}}{=}21$~mW at
  250~kbps), with compute times from measured ARM Cortex-M4 cycle counts.

  \item \textbf{Quantified energy outcome.} Direct leaf-node PQC deployment raises
  endpoint key-negotiation energy by \textbf{$+108.4\%$} over classical handshakes
  (328.46~mJ vs.\ 157.63~mJ at $N{=}100$) because the 2272-byte ML-KEM-768 handshake
  fragments into 24 link-layer frames. The proposed gateway-mediated hybrid reduces
  endpoint key-management energy by \textbf{$-96.8\%$} relative to standalone PQC and
  \textbf{$-93.3\%$} relative to classical ECDH.

  \item \textbf{Protocol-level validation.} The fragmentation model is validated in
  \emph{executed} headless Contiki-NG/COOJA runs across seven payload sizes
  (64--2272~B); the measured frame counts (1, 2, 3, 6, 11, 13, 24) match the closed form
  $\lceil(\text{payload}+12)/96\rceil$ exactly.

  \item \textbf{Formal security verification.} A suite of ProVerif~2.05 models
  machine-checks key-lifecycle secrecy and authentication, single-use refresh semantics,
  endpoint isolation under device compromise, gateway-compromise forward secrecy,
  harvest-now-decrypt-later resistance, cross-session/device key independence, ticket
  replay under packet loss, and active local-link impersonation.

  \item \textbf{Reproducibility.} Every quantitative and symbolic result is regenerable
  from a public, tagged artifact containing one run script per evidence stream,
  including the Contiki-NG patch required to transmit the 2272-byte handshake.
\end{enumerate}

\section*{Data and Reproducibility}
All simulation code, grounded parameter sets, figure-generation routines,
sensitivity-analysis scripts, ProVerif models, cycle-calibration scripts, and
Contiki-NG/COOJA scenarios are openly archived at:
\begin{center}
\url{https://github.com/prakashsurywanshi/quantum-safe-residential-iot}\\
(branch \texttt{ieee-submission}, tag \texttt{ieee-v1.0})
\end{center}
so that reviewers can independently reproduce each claim.

We believe this work offers a rigorous, hardware-grounded, and formally verified
framework for sustainable quantum-resilient IoT security, and we hope it is found
suitable for publication in the \emph{IEEE Internet of Things Journal}. We would be
glad to provide any additional information or artifacts requested.

\vspace{0.5em}
Yours sincerely,

\vspace{1.2em}
Prakash Suryawanshi\\
Qodeigence Infocircle Private Limited\\
\texttt{prakashsurywanshi1@live.com}

\end{document}
